\chapter{\texorpdfstring{INTRODUCTION}{}} %upper case only


On August 5, 2019, my husband, our one-year-old daughter, and I boarded a flight from Lahore, Pakistan, beginning a journey we believed would reshape our lives in hopeful ways. We were excited to chase our dreams, imagining a future where we would carve out our own place in the world, where we would be recognized as individuals, and our families would be known through us, rather than the other way around. What we did not fully realize then was that moving across continents also meant leaving behind the invisible roots that had always held us steady. Our daughter, Romana, and I traveled as my husband’s dependents as he pursued his Fulbright, Masters, and later Ph.D. in the US, carrying with us both hope and uncertainty, only gradually learning what it means to rebuild a sense of belonging when the support systems that once surrounded us were left behind. 

After multiple connections and nearly fifty-two hours of non-stop travel, we landed at Toledo Express Airport late that night. Waiting for us at 11:00 p.m. were Mr. George, and his daughter, and my dear friend Frances. These same people had welcomed me in August 2013, when I first arrived in Toledo as a young undergraduate exchange student, hosted by the University of Toledo through the US State Department cultural exchange program. This time, however, the situation was different. I was not a young undergraduate exchange student exploring the world on her own. This time I arrived with my family, carrying a lot of enthusiasm and hope intertwined with my new identities of a mother and spouse.

\section{Conception of the Project}

This personal experience of mine, as an immigrant parent, requires continuously navigating my place in the US as an individual, a mother, and a spouse while striving to figure out the process of negotiation for my young children growing up in a different cultural environment than where I grew up. This reality pushed me to learn the balancing act of communicating my Pakistani identity to my children while allowing them to simultaneously integrate into American society \citep{farooqPoeticAutoethnographyPath2025a}.  My attempts at informed parenting and intentional identity scaffolding led me to wonder how other Pakistani families made sense of their shifting identities, what are the challenges and stories of parenting, and what stories do children tell as they struggle to live in two sociocultural landscapes at once. How did the children feel about the approaches parents adopted? What worked and what did not? These experiences and questions transformed my lived experience into a scholarly inquiry. 

In this dissertation, I focused on a Pakistani American population based on my own positionality as I identify myself as a Sunni Muslim cisgender Pakistani mother living in the US. My lived experience and understanding of Pakistan’s cultural, linguistic, and regional diversity allowed me to interpret participants’ narratives with cultural sensitivity and contextual depth. My positionality informed my understanding of the nuanced meaning of verbal and nonverbal communication during the data collection and interpretation, which strengthened the study.

\section{Purpose of the Study}

The purpose of this study is to examine the challenges faced by first generation Pakistani immigrants (people who came to US as adults, for higher education or job and career opportunities) who are parents. The designation, one-and-a-half generation immigrants, refers to the children born in Pakistan who accompanied their parents to the US as minors.  In second-generation families, the children were born to at least one immigrant Pakistani parent in the US.  Pakistani American immigrant families must negotiate their identities both within the home and in broader social contexts. Specifically, this study investigates how parents and children engage in difficult or sensitive conversations often through storytelling at home and the ways these interactions shape identity development. 

\section{Significance}

This focus on storytelling among immigrant family members also led to my bigger contribution to the scholarship of critical interpersonal and family communication (CIFC). My positionality as a Pakistani immigrant mother and researcher empowered me to write from the perspective of an immigrant who has experienced the process of integration and adjustment of Pakistani cultural values into the American community \citep{farooqPoeticAutoethnographyPath2025a}. 

This background enabled me to be better prepared to formulate the questions and agenda that shaped the inquiry upon which this qualitative research is based. CIFC scholarship has repeatedly called for research centered on marginalized and immigrant family experiences. Responding to this call, the present study spotlights meaning-making, relational negotiation, and identity within Pakistani American families. The population of my study provides a perspective beyond the typical WEIRD (Western, Educated, Industrialized, Rich, and Democratic) sample of most studies in CIFC \citep{manningHammerCriticalTurn2021, suterTaleTwoMommies2016}. It allowed me to examine how men and women of a specific culture express themselves as belonging in their families and how these actions define their identity as individuals and family members. 

The family stories that weave together the familial and social expectations as well as values also serve as family archives \citep{koenigkellasFamilyStoriesStorytelling2012}. Furthermore, immigrant family stories carry the intense experience of migration and coexistence of contrasting identities, which create the dissonance between the cultural norms of their home of origin and those of the host country, producing identity gaps among immigrants \citep{shinCommunicationTheoryIdentity2017, jungElaboratingCommunicationTheory2004}. Immigrant family stories carry the richness of living in two worlds at once along with its price in the form of identity gaps and poor health outcomes \citep{koenigkellasCommunicatedNarrativeSenseMaking2022, shinCommunicationTheoryIdentity2017}.

The study is guided by the retrospective storytelling heuristic of Communicated Narrative Sense-Making Theory (CNSM) and the five layers of the Communication Theory of Identity (CTI) \citep{hecht2002ResearchOdyssey1993, kuiperCommunicationTheoryIdentity2021}. These paradigms illuminate how family narratives and communication processes serve as sites of identity negotiation, cultural transmission, and intergenerational meaning-making within Pakistani American families \citep{hecht2002ResearchOdyssey1993, kuiperCommunicationTheoryIdentity2021, koenigkellasCommunicatedNarrativeSenseMaking2018}. 

Marrying CNSM with CTI allows me to examine how culturally situated family stories function as a site where intergenerational meaning-making and identity negotiation unfold, extending both theories to immigrant and minority family contexts. This inquiry also addresses the whiteness and homogeneity present in studies that have tested the CNSM theory so far \citep{koenigkellasCommunicatedNarrativeSenseMaking2022}.

CTI and CNSM are particularly useful for this study because together they illuminate how identity is constructed and negotiated through storytelling. CTI provides a framework for examining identity across personal, relational, communal, and material layers \citep{kuiperCommunicationTheoryIdentity2021, hecht2002ResearchOdyssey1993, hechtCommunicationTheoryIdentity2021}, while CNSM foregrounds retrospective narratives as key sites of meaning-making within families \citep{koenigkellasCommunicatedNarrativeSenseMaking2022}. Used in combination, these theories allow the study to analyze how Pakistani families use everyday communication and storytelling to make sense of migration, parenting, and belonging across generations. 

Michael Hecht proposed the communication theory of identity (CTI) in 1993 based upon both post-positivistic and interpretive frameworks \citep{hecht2002ResearchOdyssey1993, shinCommunicationTheoryIdentity2017}. CTI draws from social identity theory and suggests that identity contains multiple frames, each representing a different aspect of one’s identity. All the frames overlap and, thus, reflect an individual's complete identity. CTI focuses on the mutual influences on both communication and identity, positing that communication is identity and that identity is a communication field \citep{hechtCommunicationTheoryIdentity2021, jungElaboratingCommunicationTheory2004, shinCommunicationTheoryIdentity2017}. Originally, \citet{hecht2002ResearchOdyssey1993} proposed the four layers of identity – personal, relational, enacted, and communal. \citet{kuiperCommunicationTheoryIdentity2021} introduced “material” as the fifth frame of identity, which refers to the physical aspects of self, health, and the environment. 

I focused on the Pakistani immigrant population in the US because of my personal identity and experiences as well as the need for stories from marginalized populations. Critical interpersonal and family communication (CIFC) scholars argue that we need research on the experiences of marginalized communities, such as immigrant families. This study presents an opportunity to explore the nuances of a different culture in the US beyond the usual emphasis on educated, prosperous people in the United States or Europe. My investigation will fill in some blanks in the literature concerning immigrant experiences as parents straddling their home and US cultural expectations. 
Insights gleaned from my role as a Pakistani immigrant parent who constantly adjusts to US community life while teaching my children about their Pakistani heritage has helped me to design this qualitative study.

\section{Positionality}

I position myself as a Pakistani Sunni Muslim Cisgender woman, born in the 1990s into a Punjabi middle-class household. I was born and brought up in Pakistan armed forces cantonments and bases in an environment marked by high academic expectations and stability. My father has served as a physician in the Pakistan Army all his life, and my mother has been a homemaker; both of my parents inherited agricultural land and a strong family network. This socioeconomic foundation provided me with consistent access to quality education, emotional support, and the resources necessary to pursue my academic and personal aspirations. 

I acknowledge my experiences may influence the analysis and interpretation of data, and I am a part of the data, which may not resonate with some. This project cannot speak for the vast and diverse Pakistani American families, particularly those shaped by different class, ethnic, religious, linguistic, or geographic contexts \citep{farooqPoeticAutoethnographyPath2025a}. Therefore, I approach this work with humility and awareness of the limitations of my perspective that is based on interviewing 24 immigrant Pakistani family members living in various cities and their suburbs in the Midwest, including Chicago, Illinois, and Toledo, Ohio. 

This study employed poetic inquiry, a qualitative, arts-based research (ABR) methodology. The methodology integrated poetic inquiry as narrative thematic analysis to systematically examine stories content, structure, and framing, providing insight into how different Pakistani families made sense of migration, parenting, and identity while attending to relational and well-being outcomes. 

Poetic inquiry is defined as “the use of poetry crafted from research endeavors, either before project analysis, as a project analysis, and/or poetry that is part of or that constitutes an entire research project” \citep[p.~210]{faulknerPoeticInquiryCraft2019}. My process involved crafting found poems that reflected respondents’ feelings, thoughts, and aspirations in a way that prose alone could not accomplish. I explain in detail how poetic inquiry conveys lived experiences distilled from interviews into meaningful, authentic poetry in the methodology section.

\section{Organization of the Chapters}

The first chapter introduces the personal, scholarly, and sociopolitical rationale for studying the communication of identity development and negotiation among Pakistani American immigrant families. The inquiry is grounded in CTI and CNSM theory that posits that the identity is an ongoing, relational, communicative process shaped by intergenerational meaning making and retrospective storytelling. This research explored the lived experiences of first-generation Pakistani immigrant families through qualitative, narrative, and art-based methodology. This study aimed to contribute insights to the field of critical interpersonal and family communication scholarships by studying a diasporic community’s family culture. 

Chapter 2, the literature review, resembles the foundation of a house. Just as builders make sure homes rise from a sturdy base, scholars do research to lay out the groundwork for their investigations. The primary intellectual bricks (books, articles, and chapters) that informed my thinking about Pakistani American families and retrospective storytelling are presented.

This groundwork naturally leads to my research questions. The full list of sources consulted appears in the bibliography at the end of this dissertation.
Chapter 3 outlines the methodology and research design, including the details of the method (i.e. in-depth episodic narrative interviews), data collection procedures, participant recruitment, and art-based, and narrative-centered thematic analytic procedures. The other topics discussed in this section include reflexivity and ethics, along with the rationale of poetic inquiry as a tool of thematic analysis. 
Chapter 4 presents the findings and analysis. I explain how the major themes emerged as found poems from the narratives of both first-generation immigrant parents and their children using the lens of CTI and CNSM. Do any of the poems arise from the one-and-half or second-generation respondents?

Chapter 5 concludes the project summarizing the findings and discussing the implications. I acknowledge the limitations of this study and offer recommendations for future research.

In summary, this study explores the personal, scholarly and sociopolitical dimensions of the communication of identity development and negotiation among Pakistani American immigrant families. The study is grounded in CTI and CNSM theory, which posits that identity is an ongoing, relational, communicative process shaped by intergenerational meaning-making and retrospective storytelling. This research explores the lived experiences of first-generation Pakistani immigrant families through qualitative, narrative and art-based methodology. This study aims to contribute meaningful insight to the field of critical interpersonal and family communication scholarships by studying diasporic community’s family culture.
